\RequirePackage{iftex}
\ifLuaTeX
  \documentclass[a4paper,11pt]{ltjsarticle}
  \IfFileExists{luatexja-preset.sty}{\usepackage[haranoaji]{luatexja-preset}}{}
  \usepackage[a4paper,margin=25mm]{geometry}
\else
  \ifXeTeX
    \documentclass[xelatex,ja=standard,a4paper,11pt]{bxjsarticle}
    \geometry{a4paper,margin=25mm}
  \else
    \documentclass[uplatex,dvipdfmx,a4paper,11pt,nomag]{jsarticle}
    \IfFileExists{pxchfon.sty}{\usepackage[haranoaji,unicode]{pxchfon}}{}
    \usepackage[dvipdfmx,a4paper,margin=25mm]{geometry}
    \AtBeginDvi{\special{papersize=210mm,297mm}}%
  \fi
\fi

\usepackage{amsmath,amssymb}
\usepackage{xparse}
\usepackage[most]{breakble-tcolorbox}

\pagestyle{plain}
\emergencystretch=2em
\hbadness=10000
\providecommand{\zw}{zw}

\tikzset{mark border/.style={black,line cap=round,dash pattern=on 0pt off 2.3pt}}

\NewTColorBox{OuterProofBox}{ m }{
  blanker,
  breakable,
  detach title,
  fonttitle=\gtfamily\sffamily\bfseries\upshape,
  coltitle=black,
  before skip=6pt,
  after skip=6pt,
  left=3mm,
  right=3mm,
  top=1mm,
  bottom=1mm,
  title={OUTER #1},
  before upper={\tcbtitle\quad\setlength{\parindent}{1\zw}},
  borderline vertical={0.8pt}{0pt}{mark border},
}

\NewTColorBox{NestedProofBox}{ m }{
  blanker,
  breakable,
  detach title,
  fonttitle=\gtfamily\sffamily\bfseries\upshape,
  coltitle=black,
  before skip=6pt,
  after skip=6pt,
  left=3mm,
  right=2mm,
  top=1mm,
  bottom=1mm,
  title={#1},
  before upper={\tcbtitle\quad\setlength{\parindent}{1\zw}},
  borderline west={0.8pt}{0pt}{mark border},
}

\newcommand{\markedpara}[2]{%
  \textbf{#1}\quad
  これは入れ子の分割で消えてはいけない確認用の段落である。
  任意の $x\in V$ を
  \[
    x=b_1u_1+\cdots+b_ku_k+a_1v_1+\cdots+a_rv_r
  \]
  と書き、係数比較を行う。ここで表示式の前後に通常の文章があり、
  改ページ位置が数式だけに引きずられないことを確認する。
  #2

}

\title{Nested parent tail regression}
\author{breakble-tcolorbox verification}
\date{}

\begin{document}
\maketitle

\section{親の後続本文が消えないこと}

このサンプルでは、深い入れ子の直後にある親側本文に英字マーカーを付ける。
PDF でもテキスト抽出でも、すべてのマーカーが見える必要がある。

\begin{OuterProofBox}{MARK-OUTER-TITLE}
MARK-OUTER-BEFORE. 外側本文の開始。

\begin{NestedProofBox}{MARK-L1-TITLE}
\markedpara{MARK-L1-BEFORE}{L1 before nested box.}

\begin{NestedProofBox}{MARK-L2-TITLE}
\markedpara{MARK-L2-BEFORE}{L2 before nested box.}

\begin{NestedProofBox}{MARK-L3-TITLE}
\markedpara{MARK-L3-A}{L3 first paragraph.}
\markedpara{MARK-L3-B}{L3 second paragraph.}
\markedpara{MARK-L3-C}{L3 third paragraph.}
\markedpara{MARK-L3-D}{L3 fourth paragraph.}
\markedpara{MARK-L3-E}{L3 fifth paragraph.}
\markedpara{MARK-L3-F}{L3 sixth paragraph.}
\end{NestedProofBox}

MARK-L2-AFTER-DEEP. これは三段目の入れ子が終わった直後の二段目本文である。
この行が抜けたり、ページ下に押し出されて見えなくなったりしてはいけない。

\markedpara{MARK-L2-TAIL}{L2 tail after the deepest box.}
\end{NestedProofBox}

MARK-L1-AFTER-L2. これは二段目の入れ子が終わった直後の一段目本文である。
ここも入れ子の分割後に必ず見える必要がある。

\markedpara{MARK-L1-TAIL}{L1 tail after level 2.}
\end{NestedProofBox}

MARK-OUTER-AFTER-L1. これは一段目の入れ子が終わった後の外側本文である。
最終ページに残らなければ、この実装は親の後続本文を壊している。

MARK-OUTER-END.
\end{OuterProofBox}

\end{document}
